% q0009.tex — The Great Bread Divide (Growth Theory / TFP)
\question{
  id        = {0009},
  title     = {The Great Bread Divide},
  source    = {OBECON},
  year      = {2026},
  round     = {Seletiva IEO -- Phase A},
  language  = {English},
  topic     = {Macro},
  subtopic  = {Growth Theory / TFP},
  type      = {objective},
  statement = {%
    On planet OBECON, there are two closed economies with no government
    intervention: Econoland and Obeland. Each economy functions efficiently,
    producing only one final good---bread---and both have many machines and
    workers devoted to increasing production. The size of the consumer market
    in both economies is also the same.

    In both Econoland and Obeland, capital and labor contribute to production
    in the same proportion, and both economies have the same total amount of
    each factor available. Based on this information, which factor explains
    why Econoland has a higher Gross Domestic Product?
  },
  choices   = {%
    \item The demand for bread is greater in Econoland than in Obeland.
    \item Econoland experiences constant returns to scale, while Obeland
          experiences increasing returns to scale.
    \item Capital is used more efficiently in Obeland than in Econoland.
    \item The level of technology and innovation in Econoland is higher than
          in Obeland.
    \item Workers are less committed in Econoland than in Obeland.
  },
  answer    = {},
  solution  = {%
    % TODO: add solution
  },
}
